\section{Быстрое преобразование Фурье}

До сих пор мы изучали комбинаторные и автоматные методы работы со строками. Однако целый пласт строковых и алгебраических задач совершенно недоступен для автоматов и хешей, но изящно решается при переходе на язык многочленов. Начнём с показательной задачи о нечётком поиске подстроки.

\subsection{Мотивирующая задача: поиск подстроки с \texorpdfstring{$\le k$}{<= k} ошибками}

Пусть дан текст $s$ длины $n$ и шаблон $t$ длины $m \le n$. Требуется найти все позиции $i \in [0, n - m]$, в которых подстрока $s[i \dots i + m - 1]$ отличается от шаблона $t$ \textbf{не более чем в $k$ позициях} (то есть расстояние Хэмминга между $s[i \dots i + m - 1]$ и $t$ не превосходит $k$).

\begin{itemize}
    \item \textbf{Случай $k = 0$ (точное совпадение):} задача решается за $O(n + m)$ с помощью префикс-функции или Z-функции.
    \item \textbf{Случай $k = 1$ (не более одной опечатки):} задачу всё ещё можно решить за линейное время $O(n + m)$ с помощью двух Z-функций. В самом деле, построим Z-функцию для строки $t\mathtt{\#}s$ "--- для каждой позиции $i$ она даст длину $z_{\text{pref}}(i)$ наибольшего общего префикса шаблона $t$ и подстроки $s[i \dots i + m - 1]$ (то есть расстояние до первой слева ошибки). Затем развернём обе строки и построим Z-функцию для $t^R\mathtt{\#}s^R$ "--- она даст длину $z_{\text{suff}}(i)$ наибольшего общего суффикса $t$ и $s[i \dots i + m - 1]$ (расстояние до первой справа ошибки). Подстрока $s[i \dots i + m - 1]$ отличается от $t$ не более чем в одной позиции тогда и только тогда, когда
    \[
        z_{\text{pref}}(i) + z_{\text{suff}}(i) \ge m - 1.
    \]
    \item \textbf{Произвольное $k$:} если разрешено хотя бы $k = 2$ или $k = 100$ ошибок, трюк с префиксами и суффиксами перестаёт работать, а бинарный поиск с хешами до очередного несовпадения потребует $O(n k \log m)$ времени, что при $k = \Theta(m)$ вырождается в квадратичное время.
\end{itemize}

Оказывается, задачу для любого $k$ можно решить за $O(|\Sigma| \cdot n \log n)$ времени, если научиться быстро перемножать многочлены.

\subsection{Умножение многочленов и общая схема БПФ}

Пусть даны два многочлена
\[
    P(x) = a_0 + a_1 x + \dots + a_{n-1} x^{n-1}, \qquad Q(x) = b_0 + b_1 x + \dots + b_{m-1} x^{m-1}.
\]
Требуется найти коэффициенты их произведения $R(x) = P(x) \cdot Q(x) = c_0 + c_1 x + \dots + c_{n+m-2} x^{n+m-2}$, которые выражаются формулой свёртки:
\[
    c_r = \sum_{i + j = r} a_i b_j.
\]
Наивное вычисление коэффициентов по этой формуле работает за $O(nm)$. Чтобы ускорить умножение до $O((n+m) \log(n+m))$, сменим способ задания многочлена: вместо вектора коэффициентов будем задавать многочлен его \textbf{значениями в наборе точек}.

\begin{proposition}
    Любой многочлен степени не выше $k - 1$ однозначно восстанавливается по своим значениям в $k$ попарно различных точках $x_0, x_1, \dots, x_{k-1}$.
\end{proposition}

\begin{proof}
    Если два многочлена $A(x)$ и $B(x)$ степени $\le k - 1$ принимают одинаковые значения в $k$ различных точках $x_0, \dots, x_{k-1}$, то их разность $D(x) = A(x) - B(x)$ имеет степень $\le k - 1$ и не менее $k$ корней, откуда по следствию из теоремы Безу $D(x) \equiv 0$. Существование же многочлена с любыми заданными значениями $y_0, \dots, y_{k-1}$ даётся явной интерполяционной формулой Лагранжа.
\end{proof}

Почему представление многочленов значениями в точках так удобно для умножения? Потому что если мы знаем значения $P(x_j)$ и $Q(x_j)$ в $k$ точках ($k > \deg P + \deg Q$), то значения их произведения $R(x) = P(x) Q(x)$ в этих же точках находятся всего за $O(k)$ умножений:
\[
    R(x_j) = P(x_j) \cdot Q(x_j), \qquad j = 0, 1, \dots, k - 1.
\]
Отсюда вытекает общая четырёхшаговая схема быстрого умножения многочленов:
\begin{enumerate}
    \item Выбрать степень двойки $k = 2^d > \deg P + \deg Q$ и дополнить коэффициенты многочленов $P$ и $Q$ старшими нулями до длины $k$. Выбрать $k$ специальных <<удобных>> точек $x_0, x_1, \dots, x_{k-1}$.
    \item \textbf{Прямое преобразование Фурье (DFT):} по коэффициентам многочленов $P$ и $Q$ вычислить их значения $P(x_j)$ и $Q(x_j)$ во всех $k$ точках за $O(k \log k)$.
    \item \textbf{Поточечное перемножение:} вычислить $R(x_j) = P(x_j) \cdot Q(x_j)$ для всех $j = 0, \dots, k - 1$ за $O(k)$.
    \item \textbf{Обратное преобразование Фурье (IDFT):} по значениям $R(x_0), \dots, R(x_{k-1})$ восстановить коэффициенты многочлена $R(x)$ за $O(k \log k)$.
\end{enumerate}

Весь секрет в том, какие именно точки $x_0, \dots, x_{k-1}$ нужно выбрать, чтобы шаги 2 и 4 выполнялись за $O(k \log k)$. В качестве таких точек выбираются \textbf{комплексные корни степени $k$ из единицы}.

\subsection{Комплексные корни из единицы и рекурсивное БПФ}

Напомним, что комплексное число на единичной окружности с аргументом $\phi$ записывается в тригонометрической и показательной форме как
\[
    e^{i\phi} = \cos \phi + i \sin \phi \in \Cm, \qquad \text{где } i^2 = -1.
\]
При перемножении двух таких чисел их аргументы складываются: $e^{i\phi_1} \cdot e^{i\phi_2} = e^{i(\phi_1 + \phi_2)}$.

Уравнение $z^k = 1$ имеет в поле комплексных чисел $\Cm$ ровно $k$ различных корней, которые делят единичную окружность на $k$ равных дуг:
\[
    \omega_j = \cos\!\left(\frac{2\pi j}{k}\right) + i \sin\!\left(\frac{2\pi j}{k}\right) = e^{2\pi i j / k} = \omega^j, \qquad j = 0, 1, \dots, k - 1,
\]
где корень $\omega = \omega_1 = e^{2\pi i / k}$ называется \textbf{первообразным (примитивным) корнем степени $k$ из единицы}, поскольку его степени $\omega^0, \omega^1, \dots, \omega^{k-1}$ порождают все $k$ корней.

\begin{center}
\begin{tikzpicture}[scale=1.9]
    \draw[-Stealth, black!60] (-1.65, 0) -- (1.75, 0) node[right, font=\footnotesize] {$\re$};
    \draw[-Stealth, black!60] (0, -1.45) -- (0, 1.45) node[above, font=\footnotesize] {$\im$};
    \draw[thick, NavyBlue!70] (0, 0) circle (1);

    \foreach \j/\label/\pos in {
        0/{$1 = \omega^0$}/above right,
        1/{$\omega^1$}/above right,
        2/{$i = \omega^2$}/above left,
        3/{$\omega^3$}/above left,
        4/{$-1 = \omega^4$}/below left,
        5/{$\omega^5$}/below left,
        6/{$-i = \omega^6$}/below right,
        7/{$\omega^7$}/below right}
    {
        \pgfmathsetmacro{\ang}{\j * 45}
        \fill[BrickRed] (\ang:1) circle (1.5pt);
        \node[\pos, font=\footnotesize] at (\ang:1.06) {\label};
        \draw[dashed, gray!60] (0,0) -- (\ang:1);
    }
\end{tikzpicture}
\par\nopagebreak\medskip
\small\textit{Рис. 7.1. Корни степени $k = 8$ из единицы на комплексной плоскости.}
\end{center}

Пусть $k = 2^d$. Покажем, как вычислить значения многочлена $P(x) = a_0 + a_1 x + \dots + a_{k-1} x^{k-1}$ во всех $k$ точках $1, \omega, \omega^2, \dots, \omega^{k-1}$ за время $O(k \log k)$.

Разделим коэффициенты многочлена $P(x)$ на две группы по чётности их индексов и составим два многочлена степени $< k / 2$:
\[
    P_0(t) = a_0 + a_2 t + a_4 t^2 + \dots + a_{k-2} t^{k/2 - 1}, \qquad
    P_1(t) = a_1 + a_3 t + a_5 t^2 + \dots + a_{k-1} t^{k/2 - 1}.
\]
Непосредственно проверяется тождество
\[
    P(x) = P_0(x^2) + x \cdot P_1(x^2).
\]
Подставим сюда $x = \omega^j$ ($j = 0, 1, \dots, k - 1$):
\begin{equation}\label{eq:fft_split}
    P(\omega^j) = P_0(\omega^{2j}) + \omega^j \cdot P_1(\omega^{2j}).
\end{equation}
Посмотрим внимательно на точки $\omega^{2j}$, в которых нам нужно вычислить многочлены половинной степени $P_0$ и $P_1$. Так как $\omega$ был первообразным корнем степени $k = 2^d$, его квадрат $\omega^2 = e^{2\pi i / (k/2)}$ является \textbf{первообразным корнем степени $k / 2 = 2^{d-1}$ из единицы}. При этом для $j$ и $j + k/2$ квадраты точек совпадают:
\[
    \omega^{2(j + k/2)} = \omega^{2j + k} = \omega^{2j} \cdot \omega^k = \omega^{2j}.
\]
Следовательно, множество точек $\{\omega^{2j}\}_{j=0}^{k-1}$ состоит ровно из $k / 2$ различных корней степени $k / 2$ из единицы: $1, \omega^2, \omega^4, \dots, \omega^{k-2}$.

Значит, чтобы найти $P(\omega^j)$ во всех $k$ точках, нам достаточно:
\begin{enumerate}
    \item рекурсивно вызвать алгоритм БПФ для многочлена $P_0$ длины $k/2$ с корнем $\omega^2$;
    \item рекурсивно вызвать алгоритм БПФ для многочлена $P_1$ длины $k/2$ с корнем $\omega^2$;
    \item за один цикл длины $k$ пересчитать значения $P(\omega^j)$ по формуле~\eqref{eq:fft_split}.
\end{enumerate}
Время работы удовлетворяет рекуррентному соотношению $T(k) = 2 T(k / 2) + O(k)$, откуда по мастер-теореме $T(k) = O(k \log k)$.

\subsection{Матрица Вандермонда и обратное БПФ}

Теперь разберёмся с шагом 4 "--- обратным преобразованием Фурье. Запишем задачу вычисления значений многочлена в точках $1, \omega, \dots, \omega^{k-1}$ на языке линейной алгебры. Пусть $a = (a_0, \dots, a_{k-1})^T$ "--- столбец коэффициентов многочлена $P(x)$, а $y = (y_0, \dots, y_{k-1})^T$ "--- столбец его значений $y_i = P(\omega^i)$. Тогда
\[
    y_i = \sum_{j=0}^{k-1} a_j \omega^{ij} \quad \Lra \quad y = W_k a,
\]
где $W_k = (\omega^{ij})_{i,j=0}^{k-1}$ "--- симметричная матрица Вандермонда размера $k \times k$:
\[
    W_k = \begin{pmatrix}
        1 & 1 & 1 & \cdots & 1 \\
        1 & \omega & \omega^2 & \cdots & \omega^{k-1} \\
        1 & \omega^2 & \omega^4 & \cdots & \omega^{2(k-1)} \\
        \vdots & \vdots & \vdots & \ddots & \vdots \\
        1 & \omega^{k-1} & \omega^{2(k-1)} & \cdots & \omega^{(k-1)^2}
    \end{pmatrix}.
\]
Задача обратного преобразования Фурье "--- выразить вектор коэффициентов $a = W_k^{-1} y$ через вектор значений $y$. Оказывается, обратная матрица $W_k^{-1}$ устроена почти так же, как и сама матрица $W_k$.

\begin{theorem}[Об обращении матрицы Фурье]\label{thm:ifft_matrix}
    Рассмотрим матрицу $V_k = (\omega^{-ij})_{i,j=0}^{k-1}$, полученную из матрицы $W_k$ заменой корня $\omega$ на обратный корень $\omega^{-1} = e^{-2\pi i / k}$. Тогда
    \[
        V_k \cdot W_k = k \cdot I_k, \qquad \text{откуда} \qquad W_k^{-1} = \frac{1}{k} V_k.
    \]
\end{theorem}

\begin{proof}
    Вычислим элемент на пересечении $p$-й строки и $q$-го столбца произведения матриц $V_k \cdot W_k$ ($0 \le p, q < k$):
    \[
        (V_k W_k)_{p, q} = \sum_{j=0}^{k-1} (V_k)_{p, j} (W_k)_{j, q} = \sum_{j=0}^{k-1} \omega^{-pj} \omega^{jq} = \sum_{j=0}^{k-1} \bigl(\omega^{q - p}\bigr)^j.
    \]
    Разберём два случая:
    \begin{itemize}
        \item \textbf{Диагональный элемент ($p = q$):} здесь $\omega^{q - p} = \omega^0 = 1$, поэтому каждое слагаемое равно $1$, и сумма равна $\sum_{j=0}^{k-1} 1 = k$.
        \item \textbf{Внедиагональный элемент ($p \ne q$):} обозначим $\varepsilon = \omega^{q - p}$. Так как $0 < |q - p| < k$, а $\omega$ "--- первообразный корень степени $k$, число $\varepsilon$ отлично от единицы ($\varepsilon \ne 1$). При этом $\varepsilon^k = (\omega^k)^{q - p} = 1^{q - p} = 1$. По формуле суммы геометрической прогрессии со знаменателем $\varepsilon \ne 1$ получаем:
        \[
            \sum_{j=0}^{k-1} \varepsilon^j = \frac{1 - \varepsilon^k}{1 - \varepsilon} = \frac{1 - 1}{1 - \varepsilon} = 0.
        \]
    \end{itemize}
    Следовательно, матрица $V_k W_k$ имеет числа $k$ на главной диагонали и нули вне неё, то есть равна $k I_k$.
\end{proof}

\begin{corollary}
    Чтобы выполнить обратное преобразование Фурье $a = \frac{1}{k} V_k y$, не нужно писать отдельный алгоритм: достаточно запустить \textbf{ту же самую функцию прямого БПФ} от вектора $y$, подставив вместо первообразного корня $\omega$ обратный корень $\omega^{-1}$ (который сам является первообразным корнем степени $k$ из единицы), а затем разделить все элементы полученного вектора на $k$.
\end{corollary}

\subsection{Нерекурсивная реализация и преобразование <<бабочки>>}

Рекурсивная реализация БПФ проста в записи, но на практике работает медленно из-за постоянного выделения памяти под временные векторы $P_0, P_1$ на каждом уровне рекурсии. Поймём, как переписать алгоритм нерекурсивно и выполнять все преобразования прямо на месте (in-place) в одном массиве длины $k = 2^d$.

Посмотрим, в каком порядке элементы исходного массива $a$ оказываются в листьях дерева рекурсии при $k = 8$ ($d = 3$):
\begin{center}
\begin{tabular}{@{}cl@{}}
    \toprule
    \textbf{Глубина рекурсии} & \textbf{Разбиение массива коэффициентов на блоки} \\
    \midrule
    Уровень $0$ & $(a_0, a_1, a_2, a_3, a_4, a_5, a_6, a_7)$ \\
    Уровень $1$ & $(a_0, a_2, a_4, a_6) \;\big|\; (a_1, a_3, a_5, a_7)$ \\
    Уровень $2$ & $(a_0, a_4) \;\big|\; (a_2, a_6) \;\big|\; (a_1, a_5) \;\big|\; (a_3, a_7)$ \\
    Уровень $3$ (листья) & $a_0 \;\big|\; a_4 \;\big|\; a_2 \;\big|\; a_6 \;\big|\; a_1 \;\big|\; a_5 \;\big|\; a_3 \;\big|\; a_7$ \\
    \bottomrule
\end{tabular}
\end{center}

Запишем индексы листьев $(0, 4, 2, 6, 1, 5, 3, 7)$ в трёхразрядной двоичной системе счисления:
\[
    000_2,\ 100_2,\ 010_2,\ 110_2,\ 001_2,\ 101_2,\ 011_2,\ 111_2.
\]
Легко видеть, что они получаются из чисел $0, 1, \dots, 7$ ($000_2$, $001_2$, $010_2$, $011_2$, $100_2$, $101_2$, $110_2$, $111_2$) \textbf{побитовым разворотом}.

\begin{proposition}
    При $k = 2^d$ на $i$-м месте нижнего уровня дерева рекурсии БПФ оказывается коэффициент $a_{\rev(i)}$, где $\rev(i)$ "--- число, полученное разворотом $d$-битной двоичной записи числа $i$ задом наперёд.
\end{proposition}

\begin{proof}
    На первом шаге рекурсии в левую половину массива попадают элементы с чётными индексами (младший бит равен $0$), а в правую "--- с нечётными индексами (младший бит равен $1$). Значит, младший бит исходного индекса определяет старший бит позиции листа. На втором шаге рекурсии аналогично второй с конца бит исходного индекса определяет второй сверху бит позиции листа, и так далее.
\end{proof}

Массив побитовых разворотов $\rev[0 \dots k-1]$ легко предподсчитать за линейное время $O(k)$ методом динамического программирования: пусть для числа $i \ge 1$ его старший единичный бит стоит в разряде $\text{oldest}$ (так что $2^{\text{oldest}} \le i < 2^{\text{oldest} + 1}$). Если отбросить этот старший бит, останется число $i - 2^{\text{oldest}} < i$, для которого $\rev$ уже вычислен, а сам бит $\text{oldest}$ после разворота $d$-битной записи перейдёт в разряд $d - 1 - \text{oldest}$:
\[
    \rev[i] = \rev\bigl[i - 2^{\text{oldest}}\bigr] + 2^{d - 1 - \text{oldest}}.
\]

Переставив элементы массива $a$ по правилу \texttt{if (i < rev[i]) swap(a[i], a[rev[i]])}, мы сразу получаем нижний уровень дерева рекурсии. Теперь будем подниматься снизу вверх, поочерёдно объединяя соседние блоки длины $s = 1, 2, 4, \dots, k/2$ в блоки удвоенной длины $2s$.

Как при объединении двух блоков длины $s$ (начиная с позиции $i$) пересчитать элементы на месте? Пусть $\varepsilon$ "--- первообразный корень степени $2s$ из единицы. Так как $\varepsilon^{2s} = 1$ и $\varepsilon$ первообразный, его половинная степень равна $\varepsilon^s = -1$. Поэтому для любого $j \in [0, s - 1]$ формула~\eqref{eq:fft_split} даёт:
\[
    \begin{cases}
        a_{\text{new}}[i + j] = a[i + j] + \varepsilon^j a[i + j + s], \\
        a_{\text{new}}[i + j + s] = a[i + j] + \varepsilon^{j + s} a[i + j + s] = a[i + j] - \varepsilon^j a[i + j + s].
    \end{cases}
\]
Положив $u = a[i + j]$ и $v = \varepsilon^j a[i + j + s]$, мы заменяем пару ячеек $(a[i + j], a[i + j + s])$ на $(u + v, u - v)$. Эта элементарная операция называется \textbf{преобразованием <<бабочки>>} (из-за формы её графа зависимостей).

\begin{center}
\begin{tikzpicture}[
    scale=1.0,
    dot/.style={circle, fill=black!80, inner sep=1.8pt},
    arr/.style={-Stealth, thick, NavyBlue}
]
    \node[dot, label=left:{$u = a[i + j]$}] (u_in) at (0, 1.2) {};
    \node[dot, label=left:{$a[i + j + s]$}] (v_in) at (0, 0) {};
    \node[dot] (v_mid) at (1.5, 0) {};

    \node[dot, label=right:{$a'[i + j] = u + v$}] (u_out) at (4.0, 1.2) {};
    \node[dot, label=right:{$a'[i + j + s] = u - v$}] (v_out) at (4.0, 0) {};

    \draw[arr] (v_in) -- node[below, font=\footnotesize] {$\times \varepsilon^j$} (v_mid);
    \node[above, font=\footnotesize, text=NavyBlue] at (1.5, 0.05) {$v$};
    \draw[arr] (u_in) -- (u_out);
    \draw[arr] (v_mid) -- (u_out);
    \draw[arr] (u_in) -- (v_out);
    \draw[arr, BrickRed] (v_mid) -- node[below, font=\footnotesize] {$-1$} (v_out);
\end{tikzpicture}
\par\nopagebreak\medskip
\small\textit{Рис. 7.2. Схема преобразования <<бабочки>>.}
\end{center}

\begin{lstlisting}[caption={Нерекурсивная реализация БПФ на месте (in-place)}]
using cd = complex<double>;
const double PI = acos(-1.0);

void fft(vector<cd>& a, bool invert) {
    int k = a.size();
    int d = __builtin_ctz(k); // $k = 2^d$

    // 1. Предподсчёт побитового разворота за O(k)
    vector<int> rev(k, 0);
    int oldest = -1;
    for (int i = 1; i < k; ++i) {
        if ((i & (i - 1)) == 0) {
            ++oldest;
        }
        rev[i] = rev[i - (1 << oldest)] + (1 << (d - 1 - oldest));
    }

    for (int i = 0; i < k; ++i) {
        if (i < rev[i]) {
            swap(a[i], a[rev[i]]);
        }
    }

    // 2. Объединение блоков снизу вверх (преобразование бабочки)
    for (int s = 1; s < k; s <<= 1) {
        double angle = PI / s * (invert ? -1 : 1);
        cd eps(cos(angle), sin(angle)); // корень степени 2s из единицы
        for (int i = 0; i < k; i += 2 * s) {
            cd w(1);
            for (int j = 0; j < s; ++j) {
                cd u = a[i + j];
                cd v = w * a[i + j + s];
                a[i + j] = u + v;
                a[i + j + s] = u - v;
                w *= eps;
            }
        }
    }

    if (invert) {
        for (cd& x : a) {
            x /= k;
        }
    }
}
\end{lstlisting}

\subsection{Точность вычислений и переход к полю \texorpdfstring{$\Z_p$}{Zp} (NTT)}

При работе с комплексными числами типа \texttt{double} неизбежно накапливается погрешность плавающей точки. Если исходные многочлены имеют целые коэффициенты, то на практике действует простое эмпирическое правило: пока точные значения коэффициентов произведения не превосходят $10^{11}\text{--}10^{12}$, стандартного типа \texttt{double} (или \texttt{long double}) с округлением \texttt{llround(a[i].real())} к ближайшему целому вполне достаточно, чтобы ответ совпал с точным без единой ошибки.

Что делать, если коэффициенты могут быть больше или если задача требует найти ответ по некоторому простому модулю $p$? Посмотрим, какие именно свойства поля $\Cm$ мы использовали при доказательстве корректности прямого и обратного БПФ:
\begin{enumerate}
    \item существование в поле элемента $\omega$ порядка $k = 2^d$ по умножению (такого, что $\omega^k = 1$, а $\omega^{k/2} = -1 \ne 1$);
    \item возможность разделить на $k$ в конце обратного БПФ (то есть обратимость элемента $k$ в поле).
\end{enumerate}

Никакая геометрия или непрерывность комплексных чисел нам на самом деле не нужна. Возьмём простое число вида
\[
    p = b \cdot 2^a + 1,
\]
где $2^a \ge k$, и будем выполнять все вычисления в конечном поле вычетов $\Z_p = \Z / p\Z$. Мультипликативная группа $\Z_p^* = \{1, 2, \dots, p - 1\}$ является циклической группой порядка $p - 1 = b \cdot 2^a$. Если $g$ "--- первообразный корень по модулю $p$ (порождающий всю группу $\Z_p^*$), то элемент $\omega = g^b \bmod p$ имеет в $\Z_p^*$ порядок ровно $2^a$.

Самый известный такой модуль в олимпиадном и алгоритмическом программировании:
\[
    p = 998\,244\,353 = 119 \cdot 2^{23} + 1.
\]
В поле $\Z_{998244353}$ существует корень $\omega = 3^{119} \bmod p = 31$ степени $2^{23} = 8\,388\,608$ из единицы, что позволяет перемножать многочлены степени до $8 \cdot 10^6$ в чистой целочисленной арифметике по модулю $p$ без вещественных чисел и ошибок округления. Такой вариант алгоритма называется \textbf{теоретико-числовым преобразованием} (англ. \textit{Number Theoretic Transform, NTT}).

\begin{note}
    Если в задаче нужно перемножить многочлены с огромными целыми коэффициентами (например, порядка $10^{25}$) \emph{без модуля}, можно выполнить NTT по трём различным простым модулям вида $b_i \cdot 2^a + 1$ (около $10^9$ каждый), а затем однозначно восстановить точные коэффициенты произведения по \textbf{Китайской теореме об остатках}.
\end{note}

\subsection{Решение задачи о поиске с \texorpdfstring{$\le k$}{<= k} несовпадениями через БПФ}

Вернёмся к нашей строковой задаче: дан текст $s$ длины $n$ и шаблон $t$ длины $m$, нужно для каждого сдвига $i \in [0, n - m]$ узнать, верно ли, что $s[i \dots i + m - 1]$ отличается от $t$ не более чем в $k$ позициях.

Переформулируем условие: расстояние Хэмминга не превосходит $k$ тогда и только тогда, когда количество \textbf{совпадающих символов} между $s[i \dots i + m - 1]$ и $t$ не меньше $m - k$.

Зафиксируем одну букву алфавита $c \in \Sigma$ и посчитаем для всех сдвигов $i$ сразу, сколько раз буква $c$ в шаблоне $t$ накладывается на такую же букву $c$ в подстроке $s[i \dots i + m - 1]$. Заведём два индикаторных массива из нулей и единиц:
\[
    S_c[u] = \begin{cases} 1, & \text{если } s[u] = c, \\ 0, & \text{иначе,} \end{cases}
    \qquad
    T_c[j] = \begin{cases} 1, & \text{если } t[j] = c, \\ 0, & \text{иначе.} \end{cases}
\]
Тогда число совпадений по букве $c$ при сдвиге $i$ равно скалярному произведению:
\[
    \text{match}_c(i) = \sum_{j=0}^{m-1} S_c[i + j] \cdot T_c[j].
\]
Эта сумма очень похожа на коэффициент произведения многочленов, с одной лишь разницей: в произведении многочленов сумма индексов перемножаемых коэффициентов постоянна, а у нас постоянна разность индексов $(i + j) - j = i$. Чтобы превратить разность в сумму, \textbf{развернём} массив шаблона $T_c$ задом наперёд и определим два многочлена:
\[
    P_c(x) = \sum_{u=0}^{n-1} S_c[u] x^u, \qquad Q_c(x) = \sum_{j=0}^{m-1} T_c[j] x^{m - 1 - j}.
\]
Перемножим многочлены $P_c(x)$ и $Q_c(x)$ с помощью БПФ за $O(n \log n)$ и посмотрим на коэффициент при мономе $x^{i + m - 1}$ в их произведении $R_c(x) = P_c(x) \cdot Q_c(x)$:
\[
    [x^{i + m - 1}] R_c(x) = \sum_{u + (m - 1 - j) = i + m - 1} S_c[u] \cdot T_c[j] = \sum_{j=0}^{m-1} S_c[i + j] \cdot T_c[j] = \text{match}_c(i).
\]
Проделав это для каждой буквы $c \in \Sigma$ и сложив количества совпадений $\sum_{c \in \Sigma} \text{match}_c(i)$, мы найдём общее число совпадений для каждого сдвига $i$ за время $O(|\Sigma| \cdot n \log n)$.

\begin{note}[Ускорение в два раза за счёт мнимой единицы]
    Поскольку комплексное БПФ всё равно работает с комплексными числами, количество умножений многочленов можно сократить вдвое "--- с $|\Sigma|$ до $\lceil |\Sigma| / 2 \rceil$.

    Разобьём буквы алфавита $\Sigma$ на пары $(c, d)$. Для обработки пары букв $(c, d)$ построим один комплексный многочлен $P(x)$ по тексту $s$ и один комплексный многочлен $Q(x)$ по развёрнутому шаблону $t$:
    \begin{itemize}
        \item в многочлене текста $P(x)$ на позиции буквы $c$ поставим $1$, на позиции буквы $d$ поставим $+i$, а на позиции всех остальных букв "--- $0$;
        \item в многочлене шаблона $Q(x)$ на позиции буквы $c$ поставим $1$, на позиции буквы $d$ поставим сопряжённое число $-i$, а на остальные позиции "--- $0$.
    \end{itemize}
    Посмотрим, какой вклад даёт перемножение символов текста и шаблона при их наложении:
    \begin{itemize}
        \item совпадение $c$ с $c$ даёт $1 \cdot 1 = 1$;
        \item совпадение $d$ с $d$ даёт $i \cdot (-i) = -i^2 = 1$;
        \item наложение $c$ на $d$ или $d$ на $c$ даёт $1 \cdot (-i) = -i$ или $i \cdot 1 = +i$ "--- \textbf{чисто мнимые числа} с нулевой вещественной частью.
    \end{itemize}
    Следовательно, вещественная часть коэффициента при мономе $x^{i + m - 1}$ в произведении $P(x) Q(x)$ в точности равна сумме совпадений $\text{match}_c(i) + \text{match}_d(i)$ сразу по двум буквам.
\end{note}
